\section{Discussion, Limitations and Future Work}
\label{sec:discussion}
\paragraph{What the study supports.}
The evaluated source representation achieves a small measured cache while
retaining all prompt positions in exact or approximate form. Its protected
regions and lossless history are concrete representation mechanisms.
Request-private reconstruction and layout-aware replay make this path
executable with complete timing and declared storage boundaries. The
complete cohort and negative contrasts are as important as the favorable
point estimates.

\paragraph{What it does not support.}
Different realized bytes prevent a byte-matched selection-policy claim.
Confidence intervals crossing zero prevent an assertion of preserved
quality relative to Top-K, VL-Cache or LOOK-M. This inspected cohort does
not establish independent confirmation. Neither the mean decode-owned
peak nor its maximum includes the entire request's live KV/activation
ownership. We explicitly defer shared-byte confirmation and exact
whole-request KV-peak certification rather than implying they have passed.

\paragraph{An architectural distinction, not a missing checkbox.}
The research repository also contains a six-component design: future-query
forecasting, sparse cross-modal evidence graphs, value-aware utility,
budgeted submodular selection, exact/prototype/residual tiers and
uncertainty-guided repair. The present storage variant is a different
execution path, not an empirical validation of that full design. Its
numerical patches are not decode-time residual repair, and its CUDA replay
graphs are not semantic evidence graphs. We withdraw those mechanisms as
explanations for these measurements. Their tests and implementation may
support engineering claims in their tested scope, but not a pretrained
quality or systems advantage without separately active experiments.

\paragraph{Baseline interpretation.}
All four compressed comparators are local adaptations; the nearest
quality-eligible published baseline has not been identified over a common
byte grid. The PrefixKV output-limit behavior and LOOK-M's large local
construction cost particularly limit generalization. Historical
paper-native PPL/ROUGE-L, MileBench and nominal-retention experiments use
different paths, models or protocols. They remain preserved separately
and are not pooled into the current COCO evidence or represented as tests
of this storage variant.

\paragraph{Statistical and workload scope.}
The cohort contains one image per request, one checkpoint and one caption
metric. Statistical adjustment controls the declared contrast families,
not bias from repeated development or a missing workload. The repeated
timing sample has six images, not 120 independent quality observations per
method. Shared-node CPU/clock variation, untested batches and the eager
execution policy constrain external validity. Captioning evidence is not
an OCR, document, video or long-context reasoning result.

\paragraph{Research directions.}
The strongest next tests are common realized-byte curves, independent
quality confirmation and complete request-memory accounting. Technical
extensions include group-aware source representations for GQA, reducing
reconstruction/launch costs on nonuniform layouts, and measuring batch
feasibility through prefill rather than decode alone. Any future prototype
or residual-repair integration must preserve original positions, numerical
contracts and complete storage accounting, and must establish an actual
quality benefit instead of relying on a reconstruction surrogate.

\paragraph{Response to methodological review.}
The accompanying response document maps the review-derived requirements
to explicit evidence and limitations: complete request timing is measured;
scoring and instrumented/unaltered output consistency are checked;
byte-matched superiority, first-token-boundary generalization, mechanism
ablations and batching remain outside the claim. A scope reduction is not
presented as experimental completion of those requirements.
